% ============================================================================
% CAPÍTULO 10: SISTEMAS POPULACIONAIS
% Bioinformática para Biologia de Sistemas
% ============================================================================

% ============================================================================
% TEMPLATE BASE PARA APRESENTAÇÕES EM BEAMER
% Bioinformática para Biologia de Sistemas
% ============================================================================

\documentclass[12pt, aspectratio=169]{beamer}

% ============================================================================
% PACOTES ESSENCIAIS
% ============================================================================
\usepackage[utf8]{inputenc}
\usepackage[T1]{fontenc}
\usepackage[brazilian]{babel}
\usepackage{amsmath, amssymb, amsthm}
\usepackage{graphicx}
\usepackage{xcolor}
\usepackage{tikz}
\usetikzlibrary{arrows.meta, shapes, positioning, calc, decorations.pathreplacing, decorations.shapes, decorations.pathmorphing}
\usepackage{listings}
\usepackage{booktabs}
\usepackage{multirow}
\usepackage{hyperref}
\usepackage{chemfig}
\usepackage{chemarr}

% ============================================================================
% CONFIGURAÇÃO DO TEMA
% ============================================================================
\usetheme{Madrid}
\usecolortheme{default}

% Cores personalizadas
\definecolor{azulescuro}{RGB}{0, 51, 102}
\definecolor{azulclaro}{RGB}{51, 153, 255}
\definecolor{verdeacido}{RGB}{102, 204, 0}
\definecolor{laranjacel}{RGB}{255, 153, 51}
\definecolor{roxodna}{RGB}{153, 51, 255}

\setbeamercolor{structure}{fg=azulescuro}
\setbeamercolor{palette primary}{bg=azulescuro,fg=white}
\setbeamercolor{palette secondary}{bg=azulclaro,fg=white}
\setbeamercolor{palette tertiary}{bg=azulescuro,fg=white}
\setbeamercolor{palette quaternary}{bg=azulclaro,fg=white}
\setbeamercolor{block title}{bg=azulescuro,fg=white}
\setbeamercolor{block body}{bg=azulclaro!10,fg=black}
\setbeamercolor{block title alerted}{bg=red!80,fg=white}
\setbeamercolor{block body alerted}{bg=red!10,fg=black}
\setbeamercolor{block title example}{bg=verdeacido!80,fg=white}
\setbeamercolor{block body example}{bg=verdeacido!10,fg=black}

% ============================================================================
% CONFIGURAÇÃO DE CÓDIGO
% ============================================================================
\lstset{
    language=Python,
    basicstyle=\ttfamily\small,
    keywordstyle=\color{azulescuro}\bfseries,
    commentstyle=\color{verdeacido},
    stringstyle=\color{laranjacel},
    numbers=left,
    numberstyle=\tiny\color{gray},
    stepnumber=1,
    numbersep=5pt,
    backgroundcolor=\color{white},
    showspaces=false,
    showstringspaces=false,
    showtabs=false,
    frame=single,
    rulecolor=\color{black},
    tabsize=2,
    captionpos=b,
    breaklines=true,
    breakatwhitespace=false,
    escapeinside={\%*}{*)},
    xleftmargin=10pt,
    xrightmargin=10pt
}

% Definir estilo para R
\lstdefinestyle{Rstyle}{
    language=R,
    basicstyle=\ttfamily\small,
    keywordstyle=\color{azulescuro}\bfseries,
    commentstyle=\color{verdeacido},
    stringstyle=\color{laranjacel}
}

% Definir estilo para MATLAB
\lstdefinestyle{MATLABstyle}{
    language=Matlab,
    basicstyle=\ttfamily\small,
    keywordstyle=\color{azulescuro}\bfseries,
    commentstyle=\color{verdeacido},
    stringstyle=\color{laranjacel}
}

% ============================================================================
% COMANDOS PERSONALIZADOS
% ============================================================================

% Comando para equações destacadas
\newcommand{\destaque}[1]{%
    \begin{alertblock}{}
        \begin{center}
            #1
        \end{center}
    \end{alertblock}
}

% Comando para definições
\newcommand{\definicao}[2]{%
    \begin{block}{Definição: #1}
        #2
    \end{block}
}

% Comando para exemplos
\newcommand{\exemplo}[2]{%
    \begin{exampleblock}{Exemplo: #1}
        #2
    \end{exampleblock}
}

% Comando para observações importantes
\newcommand{\importante}[1]{%
    \begin{alertblock}{Importante!}
        #1
    \end{alertblock}
}

% Comandos matemáticos úteis
\newcommand{\R}{\mathbb{R}}
\newcommand{\N}{\mathbb{N}}
\newcommand{\Z}{\mathbb{Z}}
\newcommand{\dif}{\mathrm{d}}
\newcommand{\parcial}[2]{\frac{\partial #1}{\partial #2}}
\newcommand{\derivada}[2]{\frac{\dif #1}{\dif #2}}
\newcommand{\vect}[1]{\boldsymbol{#1}}

% Comandos biológicos
\newcommand{\gene}[1]{\textit{#1}}
\newcommand{\proteina}[1]{\textsc{#1}}
\newcommand{\especie}[1]{\textit{#1}}

% ============================================================================
% CONFIGURAÇÃO DE RODAPÉ
% ============================================================================
\setbeamertemplate{footline}{
  \leavevmode%
  \hbox{%
  \begin{beamercolorbox}[wd=.33\paperwidth,ht=2.25ex,dp=1ex,center]{author in head/foot}%
    \usebeamerfont{author in head/foot}\insertshortauthor
  \end{beamercolorbox}%
  \begin{beamercolorbox}[wd=.34\paperwidth,ht=2.25ex,dp=1ex,center]{title in head/foot}%
    \usebeamerfont{title in head/foot}\insertshorttitle
  \end{beamercolorbox}%
  \begin{beamercolorbox}[wd=.33\paperwidth,ht=2.25ex,dp=1ex,right]{date in head/foot}%
    \usebeamerfont{date in head/foot}\insertshortdate{}\hspace*{2em}
    \insertframenumber{} / \inserttotalframenumber\hspace*{2ex}
  \end{beamercolorbox}}%
  \vskip0pt%
}

% ============================================================================
% CONFIGURAÇÃO DE NAVEGAÇÃO
% ============================================================================
\setbeamertemplate{navigation symbols}{}

% ============================================================================
% CONFIGURAÇÃO DE ITEMIZE/ENUMERATE
% ============================================================================
\setbeamertemplate{itemize items}[circle]
\setbeamertemplate{enumerate items}[default]

% ============================================================================
% FIM DO TEMPLATE
% ============================================================================


% ============================================================================
% INFORMAÇÕES DO DOCUMENTO
% ============================================================================
\title[Cap. 10: Sistemas Populacionais]{Capítulo 10: Sistemas Populacionais}
\subtitle{Bioinformática para Biologia de Sistemas}
\author{Ney Lemke}
\institute{DBF-Unesp}
\date{\today}

\begin{document}

% ============================================================================
% SLIDE DE TÍTULO
% ============================================================================
\begin{frame}
\titlepage
\end{frame}

% ============================================================================
% SUMÁRIO
% ============================================================================
\begin{frame}{Sumário}
\tableofcontents
\end{frame}

% ============================================================================
% SEÇÃO 1: INTRODUÇÃO
% ============================================================================
\section{Introdução}

\begin{frame}{Visão Geral do Capítulo}
\begin{block}{Objetivos de Aprendizagem}
\begin{itemize}
    \item Compreender modelos de dinâmica populacional
    \item Analisar interações entre espécies
    \item Estudar modelos epidemiológicos
    \item Explorar dinâmica espacial de populações
    \item Implementar simulações em Python
\end{itemize}
\end{block}

\vspace{0.3cm}

\begin{exampleblock}{Por que Estudar Sistemas Populacionais?}
\begin{itemize}
    \item Conservação de espécies ameaçadas
    \item Controle de epidemias
    \item Gestão de recursos naturais
    \item Compreensão de ecossistemas
\end{itemize}
\end{exampleblock}
\end{frame}

\begin{frame}{Dinâmica Populacional}
\definicao{Sistema Populacional}{
Conjunto de indivíduos da mesma espécie (ou espécies interagentes) cuja densidade varia no tempo e espaço devido a processos de nascimento, morte, migração e interações.
}

\vspace{0.3cm}

\begin{columns}
\column{0.5\textwidth}
\textbf{Níveis de Organização:}
\begin{itemize}
    \item \textbf{Espécie única:} crescimento, regulação
    \item \textbf{Interações:} predação, competição, mutualismo
    \item \textbf{Espacial:} metapopulações, difusão
\end{itemize}

\column{0.5\textwidth}
\textbf{Aplicações:}
\begin{itemize}
    \item Ecologia teórica
    \item Epidemiologia
    \item Biologia evolutiva
    \item Gestão ambiental
\end{itemize}
\end{columns}
\end{frame}

\begin{frame}{Processos Fundamentais}
\begin{center}
\begin{tikzpicture}[node distance=2cm, scale=0.9, transform shape]
    \node[draw, circle, fill=azulclaro!30, minimum size=2cm] (pop) {População\\$N(t)$};

    \node[above=1.5cm of pop, fill=verdeacido!30, rounded corners] (nasc) {Nascimentos: $b N$};
    \node[below=1.5cm of pop, fill=red!30, rounded corners] (morte) {Mortes: $d N$};
    \node[left=2cm of pop, fill=laranjacel!30, rounded corners] (imig) {Imigração: $I$};
    \node[right=2cm of pop, fill=roxodna!30, rounded corners] (emig) {Emigração: $E$};

    \draw[->, thick, verdeacido, line width=2pt] (nasc) -- (pop);
    \draw[->, thick, red, line width=2pt] (pop) -- (morte);
    \draw[->, thick, laranjacel, line width=2pt] (imig) -- (pop);
    \draw[->, thick, roxodna, line width=2pt] (pop) -- (emig);
\end{tikzpicture}
\end{center}

\destaque{
\begin{equation*}
\derivada{N}{t} = \underbrace{bN}_{\text{nascimento}} - \underbrace{dN}_{\text{morte}} + \underbrace{I}_{\text{imigração}} - \underbrace{E}_{\text{emigração}}
\end{equation*}
}
\end{frame}

% ============================================================================
% SEÇÃO 2: MODELOS DE CRESCIMENTO POPULACIONAL
% ============================================================================
\section{Modelos de Crescimento Populacional}

\begin{frame}{Crescimento Exponencial (Malthusiano)}
\begin{block}{Modelo Mais Simples}
População cresce proporcionalmente ao seu tamanho:
\begin{equation*}
\derivada{N}{t} = rN
\end{equation*}
onde $r = b - d$ é a taxa intrínseca de crescimento per capita.
\end{block}

\textbf{Solução analítica:}
\destaque{
\begin{equation*}
N(t) = N_0 e^{rt}
\end{equation*}
}

\begin{columns}
\column{0.5\textwidth}
\textbf{Características:}
\begin{itemize}
    \item Crescimento ilimitado ($r > 0$)
    \item Decaimento exponencial ($r < 0$)
    \item Tempo de duplicação: $t_d = \frac{\ln 2}{r}$
\end{itemize}

\column{0.5\textwidth}
\begin{center}
\begin{tikzpicture}[scale=0.7]
    \draw[->] (0,0) -- (5,0) node[right] {$t$};
    \draw[->] (0,0) -- (0,4) node[above] {$N$};
    \draw[thick, azulescuro, domain=0:4.5, smooth, samples=100]
        plot (\x, {0.5*exp(0.6*\x)});
    \node[azulescuro] at (3.5,3.5) {$r > 0$};
    \draw[thick, red, domain=0:4.5, smooth, samples=100]
        plot (\x, {3*exp(-0.4*\x)});
    \node[red] at (3,0.5) {$r < 0$};
\end{tikzpicture}
\end{center}
\end{columns}
\end{frame}

\begin{frame}{Crescimento Logístico (Verhulst)}
\begin{block}{Modelo com Capacidade de Suporte}
Incorpora limitação de recursos através da capacidade de suporte $K$:
\begin{equation*}
\derivada{N}{t} = rN\left(1 - \frac{N}{K}\right)
\end{equation*}
\end{block}

\textbf{Solução analítica:}
\begin{equation*}
N(t) = \frac{K N_0 e^{rt}}{K + N_0(e^{rt} - 1)} = \frac{K}{1 + \left(\frac{K-N_0}{N_0}\right)e^{-rt}}
\end{equation*}

\begin{columns}
\column{0.5\textwidth}
\textbf{Propriedades:}
\begin{itemize}
    \item Forma sigmoidal (curva S)
    \item $N \to K$ quando $t \to \infty$
    \item Ponto de inflexão em $N = K/2$
    \item Taxa máxima em $N = K/2$
\end{itemize}

\column{0.5\textwidth}
\begin{center}
\begin{tikzpicture}[scale=0.7]
    \draw[->] (0,0) -- (5,0) node[right] {$t$};
    \draw[->] (0,0) -- (0,4) node[above] {$N$};
    \draw[thick, azulescuro, domain=0:4.8, smooth, samples=100]
        plot (\x, {3/(1+9*exp(-1.5*\x))});
    \draw[dashed, verdeacido, thick] (0,3) -- (5,3) node[right, font=\small] {$K$};
    \draw[dashed, red] (0,1.5) -- (5,1.5) node[right, font=\tiny] {$K/2$};
    \fill[red] (1.5,1.5) circle (2pt);
\end{tikzpicture}
\end{center}
\end{columns}
\end{frame}

\begin{frame}[fragile]{Implementação: Crescimento Logístico}
\begin{lstlisting}[language=Python, caption=Modelo logístico em Python]
import numpy as np
import matplotlib.pyplot as plt
from scipy.integrate import odeint

def logistic(N, t, r, K):
    """Modelo de crescimento logistico"""
    return r * N * (1 - N/K)

# Parametros
r = 0.8    # taxa de crescimento
K = 1000   # capacidade de suporte
N0 = 50    # populacao inicial

# Tempo
t = np.linspace(0, 20, 500)

# Resolver EDO
N = odeint(logistic, N0, t, args=(r, K))

# Plotar
plt.figure(figsize=(10, 6))
plt.plot(t, N, 'b-', linewidth=2, label='N(t)')
plt.axhline(y=K, color='r', linestyle='--',
            linewidth=2, label=f'K = {K}')
plt.xlabel('Tempo')
plt.ylabel('Populacao')
plt.title('Crescimento Logistico')
plt.legend()
plt.grid(True, alpha=0.3)
\end{lstlisting}
\end{frame}

\begin{frame}{Efeito Allee}
\begin{block}{Crescimento com Densidade Crítica}
Em algumas populações, baixas densidades reduzem a taxa de crescimento (dificuldade de encontrar parceiros, defesa contra predadores):
\begin{equation*}
\derivada{N}{t} = rN\left(1 - \frac{N}{K}\right)\left(\frac{N}{A} - 1\right)
\end{equation*}
onde $A$ é a densidade crítica (limiar de Allee).
\end{block}

\begin{columns}
\column{0.5\textwidth}
\textbf{Características:}
\begin{itemize}
    \item Três equilíbrios: $N^* = 0, A, K$
    \item $N = 0$ estável, $N = A$ instável
    \item $N = K$ estável
    \item Risco de extinção se $N < A$
\end{itemize}

\column{0.5\textwidth}
\begin{center}
\begin{tikzpicture}[scale=0.7]
    \draw[->] (0,0) -- (5,0) node[right] {$N$};
    \draw[->] (0,-2) -- (0,2) node[above] {$\frac{dN}{dt}$};
    \draw[thick, azulescuro, domain=0:4.8, smooth, samples=100]
        plot (\x, {0.8*\x*(4.5-\x)*(\x-1.2)/5});
    \draw[dashed] (1.2,-2) -- (1.2,2);
    \draw[dashed] (4.5,-2) -- (4.5,2);
    \node[below, font=\tiny] at (1.2,0) {$A$};
    \node[below, font=\tiny] at (4.5,0) {$K$};
    \fill[red] (0,0) circle (2pt);
    \fill[white, draw=black] (1.2,0) circle (2pt);
    \fill[red] (4.5,0) circle (2pt);
\end{tikzpicture}
\end{center}
\textbf{Estável:} \textcolor{red}{$\bullet$} \quad
\textbf{Instável:} $\circ$
\end{columns}
\end{frame}

\begin{frame}{Modelos em Tempo Discreto}
\definicao{Equações de Diferença}{
Modelos onde o tempo avança em passos discretos: $N_{t+1} = f(N_t)$
}

\vspace{0.3cm}

\begin{columns}
\column{0.5\textwidth}
\begin{block}{Modelo de Ricker}
\begin{equation*}
N_{t+1} = N_t e^{r(1 - N_t/K)}
\end{equation*}
\end{block}

\textbf{Características:}
\begin{itemize}
    \item Pode exibir caos
    \item Usado para peixes com gerações discretas
    \item Comportamento complexo para $r$ alto
\end{itemize}

\column{0.5\textwidth}
\begin{block}{Modelo de Beverton-Holt}
\begin{equation*}
N_{t+1} = \frac{R N_t}{1 + (R-1)N_t/K}
\end{equation*}
\end{block}

\textbf{Características:}
\begin{itemize}
    \item Sempre estável
    \item $R$ = taxa de reprodução líquida
    \item Converge monotonicamente para $K$
\end{itemize}
\end{columns}
\end{frame}

\begin{frame}[fragile]{Implementação: Modelo de Ricker}
\begin{lstlisting}[language=Python, caption=Dinâmica discreta com Ricker]
import numpy as np
import matplotlib.pyplot as plt

def ricker(N, r, K):
    """Modelo de Ricker"""
    return N * np.exp(r * (1 - N/K))

# Parametros
r_values = [1.5, 2.0, 2.7]  # diferentes taxas de crescimento
K = 100
N0 = 10
t_max = 50

# Simular
fig, axes = plt.subplots(1, 3, figsize=(15, 4))

for i, r in enumerate(r_values):
    N = np.zeros(t_max)
    N[0] = N0

    for t in range(t_max - 1):
        N[t+1] = ricker(N[t], r, K)

    axes[i].plot(N, 'b.-', markersize=4)
    axes[i].axhline(y=K, color='r', linestyle='--',
                     label=f'K = {K}')
    axes[i].set_xlabel('Geracao t')
    axes[i].set_ylabel('N(t)')
    axes[i].set_title(f'r = {r}')
    axes[i].legend()
    axes[i].grid(True, alpha=0.3)

plt.tight_layout()
\end{lstlisting}
\end{frame}

\begin{frame}{Diagrama de Linha de Fase}
\begin{columns}
\column{0.5\textwidth}
\textbf{Análise Gráfica:}
\begin{itemize}
    \item Plotar $N_{t+1}$ vs $N_t$
    \item Linha de identidade: $N_{t+1} = N_t$
    \item Equilíbrios: intersecções
    \item Estabilidade: inclinação
\end{itemize}

\vspace{0.3cm}

\textbf{Critério de Estabilidade:}
\begin{equation*}
\left|\derivada{f}{N}\Bigg|_{N^*}\right| < 1
\end{equation*}

Se $> 1$: instável (oscilações ou caos)

\column{0.5\textwidth}
\begin{center}
\begin{tikzpicture}[scale=0.8]
    \draw[->] (0,0) -- (5,0) node[right] {$N_t$};
    \draw[->] (0,0) -- (0,5) node[above] {$N_{t+1}$};

    % Linha identidade
    \draw[dashed, gray] (0,0) -- (5,5);

    % Funcao de Ricker
    \draw[thick, azulescuro, domain=0:4.8, smooth, samples=100]
        plot (\x, {\x*exp(1.8*(1-\x/3.5))});

    % Equilibrio
    \fill[red] (3.5,3.5) circle (2pt);

    % Cobweb
    \draw[verdeacido, thick] (0.5,0.5) -- (0.5,1.2) -- (1.2,1.2) -- (1.2,2.3) -- (2.3,2.3) -- (2.3,3.2) -- (3.2,3.2) -- (3.2,3.5);

    \node[below] at (3.5,-0.2) {$K$};
\end{tikzpicture}
\end{center}

\textbf{Análise Cobweb:} trajetória iterativa
\end{columns}
\end{frame}

% ============================================================================
% SEÇÃO 3: INTERAÇÕES ENTRE ESPÉCIES
% ============================================================================
\section{Interações entre Espécies}

\begin{frame}{Tipos de Interações}
\begin{center}
\begin{tabular}{lccc}
\toprule
\textbf{Interação} & \textbf{Espécie 1} & \textbf{Espécie 2} & \textbf{Exemplo} \\
\midrule
Competição & $-$ & $-$ & Leões vs Hienas \\
Predação & $+$ & $-$ & Lobo vs Coelho \\
Mutualismo & $+$ & $+$ & Abelha vs Flor \\
Comensalismo & $+$ & $0$ & Rêmora vs Tubarão \\
Amensalismo & $-$ & $0$ & Árvore vs Plântula \\
Parasitismo & $+$ & $-$ & Carrapato vs Cão \\
\bottomrule
\end{tabular}
\end{center}

\vspace{0.3cm}

\begin{exampleblock}{Notação}
\begin{itemize}
    \item $(+)$ = efeito benéfico na taxa de crescimento
    \item $(-)$ = efeito prejudicial na taxa de crescimento
    \item $(0)$ = sem efeito
\end{itemize}
\end{exampleblock}
\end{frame}

\begin{frame}{Modelo Lotka-Volterra Predador-Presa}
\begin{block}{Sistema de Equações}
\begin{align*}
\derivada{V}{t} &= \alpha V - \beta V P \quad \text{(Presas)}\\
\derivada{P}{t} &= \delta \beta V P - \gamma P \quad \text{(Predadores)}
\end{align*}

\textbf{Parâmetros:}
\begin{itemize}
    \item $V$ = densidade de presas
    \item $P$ = densidade de predadores
    \item $\alpha$ = taxa de crescimento das presas
    \item $\beta$ = taxa de predação
    \item $\delta$ = eficiência de conversão (biomassa presa $\to$ predador)
    \item $\gamma$ = taxa de morte dos predadores
\end{itemize}
\end{block}
\end{frame}

\begin{frame}{Interpretação Biológica}
\begin{columns}
\column{0.5\textwidth}
\textbf{Equação das Presas:}
\begin{equation*}
\derivada{V}{t} = \underbrace{\alpha V}_{\substack{\text{crescimento}\\\text{exponencial}}} - \underbrace{\beta V P}_{\substack{\text{perda por}\\\text{predação}}}
\end{equation*}

\begin{itemize}
    \item Sem predadores: $\derivada{V}{t} = \alpha V$
    \item Taxa de encontro: lei de ação das massas ($VP$)
    \item Predação reduz crescimento
\end{itemize}

\column{0.5\textwidth}
\textbf{Equação dos Predadores:}
\begin{equation*}
\derivada{P}{t} = \underbrace{\delta \beta V P}_{\substack{\text{crescimento por}\\\text{consumo}}} - \underbrace{\gamma P}_{\substack{\text{morte}\\\text{natural}}}
\end{equation*}

\begin{itemize}
    \item Sem presas: $\derivada{P}{t} = -\gamma P$
    \item Crescimento proporcional ao consumo
    \item Morte independente de densidade
\end{itemize}
\end{columns}

\vspace{0.3cm}

\importante{
Este sistema exibe oscilações periódicas permanentes!
}
\end{frame}

\begin{frame}{Análise de Equilíbrio}
\begin{block}{Pontos de Equilíbrio}
Resolver $\derivada{V}{t} = 0$ e $\derivada{P}{t} = 0$:

\begin{enumerate}
    \item \textbf{Extinção total:} $(V^*, P^*) = (0, 0)$
    \item \textbf{Coexistência:}
    \begin{equation*}
    V^* = \frac{\gamma}{\delta\beta}, \quad P^* = \frac{\alpha}{\beta}
    \end{equation*}
\end{enumerate}
\end{block}

\textbf{Isoclinas nulas:}
\begin{itemize}
    \item \textbf{Presas:} $\derivada{V}{t} = 0 \Rightarrow V = 0$ ou $P = \alpha/\beta$
    \item \textbf{Predadores:} $\derivada{P}{t} = 0 \Rightarrow P = 0$ ou $V = \gamma/(\delta\beta)$
\end{itemize}

\begin{center}
\begin{tikzpicture}[scale=0.8]
    \draw[->] (0,0) -- (5,0) node[right] {$V$};
    \draw[->] (0,0) -- (0,4) node[above] {$P$};

    % Isoclinas
    \draw[thick, blue] (2,0) -- (2,4) node[above, font=\small] {$\frac{dP}{dt}=0$};
    \draw[thick, red] (0,2.5) -- (5,2.5) node[right, font=\small] {$\frac{dV}{dt}=0$};

    % Equilibrio
    \fill[verdeacido] (2,2.5) circle (3pt) node[above right] {$(V^*, P^*)$};

    % Trajetoria
    \draw[verdeacido, thick, ->, domain=0:360, smooth, samples=100]
        plot ({2 + 1.2*cos(\x)}, {2.5 + 0.8*sin(\x)});
\end{tikzpicture}
\end{center}
\end{frame}

\begin{frame}[fragile]{Implementação: Lotka-Volterra}
\begin{lstlisting}[language=Python, caption=Sistema predador-presa]
import numpy as np
import matplotlib.pyplot as plt
from scipy.integrate import odeint

def lotka_volterra(y, t, alpha, beta, delta, gamma):
    """Sistema Lotka-Volterra predador-presa"""
    V, P = y
    dVdt = alpha*V - beta*V*P
    dPdt = delta*beta*V*P - gamma*P
    return [dVdt, dPdt]

# Parametros
alpha = 1.0    # crescimento presas
beta = 0.1     # taxa de predacao
delta = 0.75   # eficiencia conversao
gamma = 1.5    # morte predadores

# Condicoes iniciais
y0 = [40, 9]   # [V0, P0]
t = np.linspace(0, 50, 1000)

# Resolver
sol = odeint(lotka_volterra, y0, t,
             args=(alpha, beta, delta, gamma))

# Plotar series temporais
plt.figure(figsize=(12, 5))
plt.subplot(1, 2, 1)
plt.plot(t, sol[:, 0], 'b-', linewidth=2, label='Presas (V)')
plt.plot(t, sol[:, 1], 'r-', linewidth=2, label='Predadores (P)')
plt.xlabel('Tempo')
plt.ylabel('Densidade')
plt.legend()
plt.grid(True, alpha=0.3)
\end{lstlisting}
\end{frame}

\begin{frame}[fragile]{Retrato de Fase}
\begin{columns}
\column{0.45\textwidth}
\begin{lstlisting}[language=Python, caption=Plano de fase, basicstyle=\ttfamily\tiny]
# Retrato de fase
plt.subplot(1, 2, 2)
plt.plot(sol[:, 0], sol[:, 1],
         'g-', linewidth=2)
plt.plot(sol[0, 0], sol[0, 1],
         'go', markersize=10,
         label='Inicio')

# Equilibrio
V_eq = gamma/(delta*beta)
P_eq = alpha/beta
plt.plot(V_eq, P_eq, 'r*',
         markersize=15,
         label='Equilibrio')

# Isoclinas
V_range = np.linspace(0, 60, 100)
plt.axhline(y=P_eq, color='r',
            linestyle='--',
            alpha=0.5,
            label='dV/dt=0')
plt.axvline(x=V_eq, color='b',
            linestyle='--',
            alpha=0.5,
            label='dP/dt=0')

plt.xlabel('Presas (V)')
plt.ylabel('Predadores (P)')
plt.legend(fontsize=8)
plt.grid(True, alpha=0.3)
plt.tight_layout()
\end{lstlisting}

\column{0.55\textwidth}
\begin{center}
\begin{tikzpicture}[scale=0.9]
    \draw[->] (0,0) -- (6,0) node[right] {Presas ($V$)};
    \draw[->] (0,0) -- (0,5) node[above] {Predadores ($P$)};

    % Isoclinas
    \draw[dashed, red, thick] (0,2.5) -- (6,2.5);
    \draw[dashed, blue, thick] (3,0) -- (3,5);

    % Trajetorias
    \draw[verdeacido, very thick, ->, domain=0:350, smooth, samples=100]
        plot ({3 + 2*cos(\x)}, {2.5 + 1.5*sin(\x)});
    \draw[verdeacido, thick, ->, domain=0:350, smooth, samples=100]
        plot ({3 + 1*cos(\x)}, {2.5 + 0.75*sin(\x)});

    % Equilibrio
    \fill[red] (3,2.5) circle (3pt) node[above right, font=\small] {$(V^*, P^*)$};

    % Campo vetorial (simplificado)
    \foreach \x in {1,2,4,5} {
        \foreach \y in {1,2,3,4} {
            \draw[->, gray, thin] (\x,\y) -- +({0.15*(\x - 0.1*\x*\y)}, {0.15*(0.05*\x*\y - \y)});
        }
    }
\end{tikzpicture}
\end{center}

\textbf{Observações:}
\begin{itemize}
    \item Órbitas fechadas (ciclos)
    \item Período depende de CI
    \item Centro neutro (não atrator)
\end{itemize}
\end{columns}
\end{frame}

\begin{frame}{Competição: Lotka-Volterra}
\begin{block}{Sistema de Competição}
Duas espécies competindo pelos mesmos recursos:
\begin{align*}
\derivada{N_1}{t} &= r_1 N_1\left(1 - \frac{N_1 + \alpha_{12}N_2}{K_1}\right)\\
\derivada{N_2}{t} &= r_2 N_2\left(1 - \frac{N_2 + \alpha_{21}N_1}{K_2}\right)
\end{align*}

\textbf{Coeficientes de competição:}
\begin{itemize}
    \item $\alpha_{12}$ = efeito da espécie 2 sobre a espécie 1
    \item $\alpha_{21}$ = efeito da espécie 1 sobre a espécie 2
\end{itemize}
\end{block}

\textbf{Interpretação:} $\alpha_{12}N_2$ converte indivíduos da espécie 2 em "equivalentes" da espécie 1 em termos de uso de recursos.
\end{frame}

\begin{frame}{Exclusão Competitiva}
\begin{block}{Quatro Cenários Possíveis}
\begin{enumerate}
    \item \textbf{Espécie 1 vence:} $K_1 > \alpha_{12}K_2$ e $K_2 < \alpha_{21}K_1$
    \item \textbf{Espécie 2 vence:} $K_2 > \alpha_{21}K_1$ e $K_1 < \alpha_{12}K_2$
    \item \textbf{Coexistência estável:} $K_1 > \alpha_{12}K_2$ e $K_2 > \alpha_{21}K_1$
    \item \textbf{Exclusão dependente de CI:} $K_1 < \alpha_{12}K_2$ e $K_2 < \alpha_{21}K_1$
\end{enumerate}
\end{block}

\begin{center}
\begin{tikzpicture}[scale=0.7]
    % Coexistencia
    \begin{scope}
        \draw[->] (0,0) -- (3,0) node[right, font=\tiny] {$N_1$};
        \draw[->] (0,0) -- (0,3) node[above, font=\tiny] {$N_2$};
        \draw[thick, blue] (2.5,0) -- (0,2);
        \draw[thick, red] (2,0) -- (0,2.5);
        \fill[verdeacido] (1.2,1.2) circle (2pt);
        \node[below, font=\tiny] at (1.5,-0.3) {Coexistência};
    \end{scope}

    % Sp 1 vence
    \begin{scope}[xshift=4cm]
        \draw[->] (0,0) -- (3,0) node[right, font=\tiny] {$N_1$};
        \draw[->] (0,0) -- (0,3) node[above, font=\tiny] {$N_2$};
        \draw[thick, blue] (2.5,0) -- (0,1.5);
        \draw[thick, red] (1.5,0) -- (0,2.5);
        \fill[verdeacido] (2.5,0) circle (2pt);
        \node[below, font=\tiny] at (1.5,-0.3) {Sp. 1 vence};
    \end{scope}
\end{tikzpicture}
\end{center}

\importante{
Princípio da Exclusão Competitiva (Gause): duas espécies não podem coexistir no mesmo nicho ecológico.
}
\end{frame}

\begin{frame}{Mutualismo}
\begin{block}{Modelo de Mutualismo}
Interação benéfica mútua (polinizadores-plantas, micorrizas-raízes):
\begin{align*}
\derivada{N_1}{t} &= r_1 N_1\left(1 - \frac{N_1}{K_1 + \alpha_{12}N_2}\right)\\
\derivada{N_2}{t} &= r_2 N_2\left(1 - \frac{N_2}{K_2 + \alpha_{21}N_1}\right)
\end{align*}
\end{block}

\textbf{Características:}
\begin{itemize}
    \item Coeficientes $\alpha_{ij} > 0$ aumentam capacidade de suporte
    \item Mutualismo facultativo: espécies podem sobreviver sozinhas
    \item Mutualismo obrigatório: requer ambas espécies
\end{itemize}

\exemplo{Abelhas e Flores}{
Abelhas aumentam reprodução das plantas (polinização)\\
Flores fornecem néctar/pólen para abelhas\\
Coevolução: dependência mútua crescente
}
\end{frame}

\begin{frame}{Modelo Nicholson-Bailey}
\begin{block}{Hospedeiro-Parasitóide em Tempo Discreto}
\begin{align*}
H_{t+1} &= \lambda H_t e^{-aP_t}\\
P_{t+1} &= c H_t (1 - e^{-aP_t})
\end{align*}

\textbf{Parâmetros:}
\begin{itemize}
    \item $H_t$ = densidade de hospedeiros na geração $t$
    \item $P_t$ = densidade de parasitóides
    \item $\lambda$ = taxa de reprodução do hospedeiro
    \item $a$ = taxa de ataque do parasitóide
    \item $c$ = número de parasitóides emergindo por hospedeiro
\end{itemize}
\end{block}

\textbf{Distribuição de Poisson:} $e^{-aP_t}$ é a probabilidade de um hospedeiro escapar do parasitismo.

\begin{exampleblock}{Aplicação}
Controle biológico de pragas agrícolas
\end{exampleblock}
\end{frame}

% ============================================================================
% SEÇÃO 4: MODELOS EPIDEMIOLÓGICOS
% ============================================================================
\section{Modelos Epidemiológicos}

\begin{frame}{Epidemiologia Matemática}
\definicao{Modelo Epidemiológico}{
Modelo matemático que descreve a transmissão e propagação de doenças infecciosas em uma população.
}

\vspace{0.3cm}

\begin{columns}
\column{0.5\textwidth}
\textbf{Compartimentos:}
\begin{itemize}
    \item \textbf{S:} Susceptíveis
    \item \textbf{E:} Expostos (incubação)
    \item \textbf{I:} Infectados (infecciosos)
    \item \textbf{R:} Recuperados (imunes)
    \item \textbf{D:} Mortos
\end{itemize}

\column{0.5\textwidth}
\textbf{Aplicações:}
\begin{itemize}
    \item Previsão de epidemias
    \item Avaliação de intervenções
    \item Políticas de vacinação
    \item Alocação de recursos
\end{itemize}
\end{columns}

\vspace{0.3cm}

\begin{exampleblock}{Exemplos Históricos}
COVID-19, Influenza, Ebola, HIV/AIDS, Malária, Dengue
\end{exampleblock}
\end{frame}

\begin{frame}{Modelo SIR}
\begin{block}{Sistema SIR Clássico}
\begin{align*}
\derivada{S}{t} &= -\beta \frac{SI}{N}\\
\derivada{I}{t} &= \beta \frac{SI}{N} - \gamma I\\
\derivada{R}{t} &= \gamma I
\end{align*}

\textbf{Parâmetros:}
\begin{itemize}
    \item $\beta$ = taxa de transmissão (contatos infectantes por unidade de tempo)
    \item $\gamma$ = taxa de recuperação ($1/\gamma$ = período infeccioso)
    \item $N = S + I + R$ = população total (constante)
\end{itemize}
\end{block}

\begin{center}
\begin{tikzpicture}[node distance=2.5cm, auto,
    box/.style={rectangle, draw, fill=azulclaro!20, minimum width=1.5cm, minimum height=1cm, rounded corners}]

    \node[box, fill=verdeacido!20] (S) {$S$};
    \node[box, fill=red!20, right of=S] (I) {$I$};
    \node[box, fill=azulescuro!20, right of=I] (R) {$R$};

    \draw[->, thick, line width=2pt] (S) -- node[above] {$\beta SI/N$} (I);
    \draw[->, thick, line width=2pt] (I) -- node[above] {$\gamma I$} (R);
\end{tikzpicture}
\end{center}
\end{frame}

\begin{frame}{Número Básico de Reprodução}
\definicao{$R_0$ (R-naught)}{
Número médio de infecções secundárias causadas por um indivíduo infectado introduzido em uma população completamente susceptível.
}

\destaque{
\begin{equation*}
R_0 = \frac{\beta}{\gamma} = \beta \times \frac{1}{\gamma} = \text{(taxa de contato)} \times \text{(duração da infecção)}
\end{equation*}
}

\begin{block}{Interpretação Epidemiológica}
\begin{itemize}
    \item $R_0 < 1$: epidemia se extingue (cada infectado infecta $< 1$ pessoa)
    \item $R_0 = 1$: limiar epidêmico (endemia estável)
    \item $R_0 > 1$: epidemia se propaga
\end{itemize}
\end{block}

\begin{exampleblock}{Valores Típicos de $R_0$}
\begin{itemize}
    \item Sarampo: 12-18 (altamente contagioso)
    \item COVID-19 (original): 2-3
    \item Influenza sazonal: 1-2
    \item Ebola: 1.5-2.5
\end{itemize}
\end{exampleblock}
\end{frame}

\begin{frame}{Limiar de Vacinação}
\begin{block}{Fração Crítica de Vacinação}
Para eliminar a doença, é necessário vacinar uma fração $f$ da população:
\begin{equation*}
f > 1 - \frac{1}{R_0}
\end{equation*}
\end{block}

\textbf{Derivação:} Vacinação reduz susceptíveis efetivos. Para controle, precisamos $R_0^{eff} = R_0(1-f) < 1$.

\begin{center}
\begin{tikzpicture}[scale=0.9]
    \draw[->] (0,0) -- (7,0) node[right] {$R_0$};
    \draw[->] (0,0) -- (0,4) node[above] {$f$ (\%)};

    \draw[thick, azulescuro, domain=1:6.5, smooth, samples=100]
        plot (\x, {3*(1 - 1/\x)});

    \draw[dashed] (1,0) -- (1,4);
    \node[below] at (1,0) {1};

    \foreach \x/\label in {2/50\%, 3/67\%, 4/75\%, 5/80\%} {
        \fill[red] (\x, {3*(1-1/\x)}) circle (2pt);
        \node[right, font=\tiny] at (\x, {3*(1-1/\x)}) {\label};
    }

    \node[fill=yellow!30, rounded corners, font=\small] at (4,0.8) {Zona de controle};
\end{tikzpicture}
\end{center}

\exemplo{Sarampo ($R_0 \approx 15$)}{
$f > 1 - 1/15 = 0.93$ (93\% de cobertura vacinal necessária)
}
\end{frame}

\begin{frame}[fragile]{Implementação: Modelo SIR}
\begin{lstlisting}[language=Python, caption=Simulação de epidemia SIR]
import numpy as np
import matplotlib.pyplot as plt
from scipy.integrate import odeint

def sir_model(y, t, beta, gamma, N):
    """Modelo SIR"""
    S, I, R = y
    dSdt = -beta * S * I / N
    dIdt = beta * S * I / N - gamma * I
    dRdt = gamma * I
    return [dSdt, dIdt, dRdt]

# Parametros
N = 10000        # populacao total
beta = 0.5       # taxa de transmissao (1/dia)
gamma = 0.1      # taxa de recuperacao (1/dia)
R0 = beta/gamma  # R0 = 5

# Condicoes iniciais
I0 = 10
S0 = N - I0
R0_initial = 0
y0 = [S0, I0, R0_initial]

# Tempo
t = np.linspace(0, 200, 1000)

# Resolver
sol = odeint(sir_model, y0, t, args=(beta, gamma, N))
S, I, R = sol.T
\end{lstlisting}
\end{frame}

\begin{frame}[fragile]{Análise da Epidemia}
\begin{lstlisting}[language=Python, caption=Análise de resultados]
# Plotar
plt.figure(figsize=(12, 5))

plt.subplot(1, 2, 1)
plt.plot(t, S, 'b-', linewidth=2, label='Suscetiveis (S)')
plt.plot(t, I, 'r-', linewidth=2, label='Infectados (I)')
plt.plot(t, R, 'g-', linewidth=2, label='Recuperados (R)')
plt.xlabel('Tempo (dias)')
plt.ylabel('Numero de individuos')
plt.title(f'Dinamica SIR ($R_0$ = {R0:.1f})')
plt.legend()
plt.grid(True, alpha=0.3)

# Encontrar pico da epidemia
peak_idx = np.argmax(I)
peak_time = t[peak_idx]
peak_infected = I[peak_idx]

print(f"Pico de infectados: {peak_infected:.0f} individuos")
print(f"Tempo do pico: {peak_time:.1f} dias")
print(f"Taxa de ataque final: {R[-1]/N*100:.1f}%")

# Limiar epidemico
S_threshold = N / R0
print(f"Limiar de suscetiveis: {S_threshold:.0f}")

plt.subplot(1, 2, 2)
plt.plot(S, I, 'purple', linewidth=2)
plt.axvline(x=S_threshold, color='red', linestyle='--',
            label=f'S = N/R0 = {S_threshold:.0f}')
plt.xlabel('Suscetiveis (S)')
plt.ylabel('Infectados (I)')
plt.title('Retrato de Fase')
plt.legend()
plt.grid(True, alpha=0.3)
\end{lstlisting}
\end{frame}

\begin{frame}{Modelo SIS}
\begin{block}{SIS: Sem Imunidade Permanente}
Indivíduos recuperados voltam a ser susceptíveis (ex: resfriado comum, gonorreia):
\begin{align*}
\derivada{S}{t} &= -\beta \frac{SI}{N} + \gamma I\\
\derivada{I}{t} &= \beta \frac{SI}{N} - \gamma I
\end{align*}
\end{block}

\begin{center}
\begin{tikzpicture}[node distance=3cm, auto,
    box/.style={rectangle, draw, fill=azulclaro!20, minimum width=1.5cm, minimum height=1cm, rounded corners}]

    \node[box, fill=verdeacido!20] (S) {$S$};
    \node[box, fill=red!20, right of=S] (I) {$I$};

    \draw[->, thick, line width=2pt, bend left=20] (S) to node[above] {$\beta SI/N$} (I);
    \draw[->, thick, line width=2pt, bend left=20] (I) to node[below] {$\gamma I$} (S);
\end{tikzpicture}
\end{center}

\textbf{Equilíbrio endêmico:} Se $R_0 > 1$,
\begin{equation*}
I^* = N\left(1 - \frac{1}{R_0}\right), \quad S^* = \frac{N}{R_0}
\end{equation*}

\importante{
Doença persiste permanentemente na população (endemia)!
}
\end{frame}

\begin{frame}{Modelo SEIR}
\begin{block}{SEIR: Com Período de Incubação}
Adiciona compartimento de Expostos (infectados mas não infecciosos):
\begin{align*}
\derivada{S}{t} &= -\beta \frac{SI}{N}\\
\derivada{E}{t} &= \beta \frac{SI}{N} - \sigma E\\
\derivada{I}{t} &= \sigma E - \gamma I\\
\derivada{R}{t} &= \gamma I
\end{align*}

onde $\sigma$ = taxa de progressão ($1/\sigma$ = período de incubação).
\end{block}

\begin{center}
\begin{tikzpicture}[node distance=2cm, auto,
    box/.style={rectangle, draw, fill=azulclaro!20, minimum width=1.2cm, minimum height=0.8cm, rounded corners, font=\small}]

    \node[box, fill=verdeacido!20] (S) {$S$};
    \node[box, fill=yellow!30, right of=S] (E) {$E$};
    \node[box, fill=red!20, right of=E] (I) {$I$};
    \node[box, fill=azulescuro!20, right of=I] (R) {$R$};

    \draw[->, thick] (S) -- node[above, font=\tiny] {$\beta SI/N$} (E);
    \draw[->, thick] (E) -- node[above, font=\tiny] {$\sigma E$} (I);
    \draw[->, thick] (I) -- node[above, font=\tiny] {$\gamma I$} (R);
\end{tikzpicture}
\end{center}

\textbf{$R_0$ modificado:}
\begin{equation*}
R_0 = \frac{\beta}{\gamma} \quad \text{(permanece o mesmo)}
\end{equation*}

mas a dinâmica temporal é alterada (pico mais tardio).
\end{frame}

\begin{frame}{Extensões do Modelo SIR}
\begin{columns}
\column{0.5\textwidth}
\textbf{Demografia:}
\begin{itemize}
    \item Nascimentos: $\mu N$
    \item Mortes naturais: $\mu S, \mu I, \mu R$
    \item Mortalidade pela doença: $\alpha I$
\end{itemize}

\textbf{Estrutura etária:}
\begin{itemize}
    \item Diferentes $\beta$ por faixa etária
    \item Matrizes de contato
    \item Vacinação por idade
\end{itemize}

\column{0.5\textwidth}
\textbf{Intervenções:}
\begin{itemize}
    \item Vacinação: $\nu S$
    \item Quarentena: redução de $\beta$
    \item Tratamento: aumento de $\gamma$
    \item Distanciamento social
\end{itemize}

\textbf{Espaço:}
\begin{itemize}
    \item Modelos metapopulacionais
    \item Difusão espacial
    \item Redes de contato
\end{itemize}
\end{columns}

\vspace{0.3cm}

\begin{exampleblock}{Modelagem de COVID-19}
Combinou SEIR + demografia + estrutura etária + intervenções não-farmacêuticas + vacinação
\end{exampleblock}
\end{frame}

% ============================================================================
% SEÇÃO 5: DINÂMICA ESPACIAL
% ============================================================================
\section{Dinâmica Espacial}

\begin{frame}{Metapopulações}
\definicao{Metapopulação}{
Conjunto de populações locais conectadas por dispersão, onde cada população local pode sofrer extinção e recolonização.
}

\vspace{0.3cm}

\begin{columns}
\column{0.5\textwidth}
\textbf{Modelo de Levins:}
\begin{equation*}
\derivada{p}{t} = cp(1-p) - ep
\end{equation*}

\begin{itemize}
    \item $p$ = fração de patches ocupados
    \item $c$ = taxa de colonização
    \item $e$ = taxa de extinção local
\end{itemize}

\textbf{Equilíbrio:}
\begin{equation*}
p^* = 1 - \frac{e}{c}
\end{equation*}

\column{0.5\textwidth}
\begin{center}
\begin{tikzpicture}[scale=0.8]
    % Patches
    \foreach \x/\y/\occupied in {0/0/1, 2/0/0, 1/1.5/1, 3/1.5/1, 0.5/3/0, 2.5/3/1} {
        \ifnum\occupied=1
            \fill[verdeacido!50] (\x,\y) circle (0.4);
        \else
            \draw[thick, gray] (\x,\y) circle (0.4);
        \fi
    }

    % Dispersao (setas)
    \draw[->, thick, azulescuro] (0.3,0.3) -- (0.8,1.2);
    \draw[->, thick, azulescuro] (1.2,1.8) -- (2.2,1.8);
    \draw[->, thick, azulescuro] (3,2) -- (2.6,2.7);

    \node[below, font=\small] at (1.5,-0.5) {Ocupados: verde};
    \node[below, font=\small] at (1.5,-1) {Vazios: cinza};
\end{tikzpicture}
\end{center}

\textbf{Condição de persistência:}
\begin{equation*}
c > e
\end{equation*}
\end{columns}
\end{frame}

\begin{frame}{Equações de Reação-Difusão}
\begin{block}{Difusão Espacial}
Combina crescimento local com dispersão espacial:
\begin{equation*}
\parcial{N}{t} = \underbrace{f(N)}_{\text{reação}} + \underbrace{D\nabla^2 N}_{\text{difusão}}
\end{equation*}

Em 1D:
\begin{equation*}
\parcial{N}{t} = f(N) + D\frac{\partial^2 N}{\partial x^2}
\end{equation*}
\end{block}

\textbf{Parâmetros:}
\begin{itemize}
    \item $D$ = coeficiente de difusão (taxa de dispersão)
    \item $f(N)$ = dinâmica local (ex: logística)
    \item $\nabla^2$ = operador Laplaciano
\end{itemize}

\exemplo{Equação de Fisher-KPP}{
\begin{equation*}
\parcial{N}{t} = rN\left(1 - \frac{N}{K}\right) + D\frac{\partial^2 N}{\partial x^2}
\end{equation*}
Modelo de invasão espacial de uma espécie
}
\end{frame}

\begin{frame}{Ondas Viajantes}
\begin{block}{Solução de Onda Viajante}
Perturbação que se propaga no espaço mantendo sua forma:
\begin{equation*}
N(x,t) = N(z), \quad z = x - ct
\end{equation*}
onde $c$ é a velocidade da onda.
\end{block}

\textbf{Velocidade mínima (Fisher-KPP):}
\begin{equation*}
c_{min} = 2\sqrt{rD}
\end{equation*}

\begin{center}
\begin{tikzpicture}[scale=0.8]
    \draw[->] (0,0) -- (8,0) node[right] {$x$};
    \draw[->] (0,0) -- (0,3) node[above] {$N$};

    % Onda em t1
    \draw[thick, blue, domain=0:7, smooth, samples=100]
        plot (\x, {2.5/(1+exp(-2*(\x-2)))});
    \node[blue] at (1,2.5) {$t_1$};

    % Onda em t2
    \draw[thick, red, domain=0:7, smooth, samples=100]
        plot (\x, {2.5/(1+exp(-2*(\x-4)))});
    \node[red] at (3,2.5) {$t_2$};

    % Velocidade
    \draw[<->, thick, verdeacido] (2,0.3) -- (4,0.3) node[midway, below] {$c \Delta t$};

    \draw[dashed] (0,2.5) -- (8,2.5) node[right, font=\tiny] {$K$};
\end{tikzpicture}
\end{center}

\textbf{Aplicações:} Invasões biológicas, frentes de epidemias, propagação de genes
\end{frame}

\begin{frame}{Padrões de Turing}
\begin{block}{Instabilidade de Difusão}
Sistema de duas espécies com difusão diferencial pode gerar padrões espaciais espontâneos:
\begin{align*}
\parcial{u}{t} &= f(u,v) + D_u \nabla^2 u\\
\parcial{v}{t} &= g(u,v) + D_v \nabla^2 v
\end{align*}

\textbf{Condição para padrões:} $D_v \gg D_u$ (inibidor difunde mais rápido que ativador)
\end{block}

\begin{columns}
\column{0.5\textwidth}
\textbf{Mecanismo:}
\begin{itemize}
    \item Ativador de curto alcance
    \item Inibidor de longo alcance
    \item Desestabilização por difusão
\end{itemize}

\textbf{Exemplos biológicos:}
\begin{itemize}
    \item Padrões de pelagem animal
    \item Distribuição de estômatos em folhas
    \item Arranjo de penas em aves
\end{itemize}

\column{0.5\textwidth}
\begin{center}
\begin{tikzpicture}[scale=0.7]
    % Simulacao simplificada de padroes
    \foreach \x in {0,1,...,6} {
        \foreach \y in {0,1,...,4} {
            \pgfmathsetmacro{\val}{sin(60*\x + 45)*sin(60*\y + 30)}
            \ifnum\val>0
                \fill[azulescuro, opacity=\val] (\x*0.6, \y*0.6) circle (0.25);
            \else
                \fill[laranjacel, opacity={-\val}] (\x*0.6, \y*0.6) circle (0.25);
            \fi
        }
    }
    \node[below, font=\small] at (2,-0.5) {Padrão de manchas};
\end{tikzpicture}
\end{center}
\end{columns}

\importante{
Alan Turing (1952): "The Chemical Basis of Morphogenesis"
}
\end{frame}

\begin{frame}[fragile]{Implementação: Difusão 1D}
\begin{lstlisting}[language=Python, caption=Equação de Fisher-KPP em 1D, basicstyle=\ttfamily\footnotesize]
import numpy as np
import matplotlib.pyplot as plt

def fisher_kpp_1d(N, t, r, K, D, dx):
    """Equacao de Fisher-KPP em 1D"""
    # Termo de reacao (logistico)
    reaction = r * N * (1 - N/K)

    # Termo de difusao (diferenca finita central)
    diffusion = np.zeros_like(N)
    diffusion[1:-1] = D * (N[2:] - 2*N[1:-1] + N[:-2]) / dx**2

    # Condicoes de contorno (sem fluxo)
    diffusion[0] = diffusion[-1] = 0

    return reaction + diffusion

# Parametros
L = 50          # tamanho do dominio
nx = 200        # pontos espaciais
x = np.linspace(0, L, nx)
dx = x[1] - x[0]

r = 0.5         # taxa de crescimento
K = 1.0         # capacidade de suporte
D = 0.1         # coeficiente de difusao

# Condicao inicial (pulso localizado)
N0 = np.zeros(nx)
N0[nx//4:nx//4+10] = K

# Simular usando metodo de Euler
t_max = 100
dt = 0.01
steps = int(t_max / dt)

N = N0.copy()
for step in range(steps):
    N += dt * fisher_kpp_1d(N, step*dt, r, K, D, dx)
\end{lstlisting}
\end{frame}

\begin{frame}{Metapopulações e Redes}
\begin{columns}
\column{0.5\textwidth}
\textbf{Modelo em Rede:}
\begin{equation*}
\derivada{N_i}{t} = f(N_i) + \sum_{j=1}^{n} m_{ij}(N_j - N_i)
\end{equation*}

\begin{itemize}
    \item $N_i$ = população no patch $i$
    \item $m_{ij}$ = taxa de migração de $i$ para $j$
    \item $f(N_i)$ = dinâmica local
\end{itemize}

\textbf{Sincronização:}
\begin{itemize}
    \item Alta conectividade $\to$ sincronismo
    \item Baixa conectividade $\to$ assincronia
    \item Importante para estabilidade
\end{itemize}

\column{0.5\textwidth}
\begin{center}
\begin{tikzpicture}[scale=0.9,
    patch/.style={circle, draw, fill=azulclaro!30, minimum size=0.8cm, font=\small}]

    \node[patch] (1) at (0,2) {$N_1$};
    \node[patch] (2) at (2,2) {$N_2$};
    \node[patch] (3) at (1,0) {$N_3$};
    \node[patch] (4) at (3,0) {$N_4$};
    \node[patch] (5) at (4,2) {$N_5$};

    \draw[<->, thick] (1) -- (2);
    \draw[<->, thick] (2) -- (3);
    \draw[<->, thick] (3) -- (4);
    \draw[<->, thick] (2) -- (4);
    \draw[<->, thick] (2) -- (5);
    \draw[<->, thick] (4) -- (5);

    \node[below, font=\small] at (2,-1) {Rede de patches};
\end{tikzpicture}
\end{center}

\textbf{Topologias:}
\begin{itemize}
    \item Regular (grade)
    \item Aleatória
    \item Escala-livre
    \item Mundo-pequeno
\end{itemize}
\end{columns}

\begin{exampleblock}{Aplicação}
Conservação: corredores ecológicos conectando fragmentos de habitat
\end{exampleblock}
\end{frame}

% ============================================================================
% SEÇÃO 6: EXERCÍCIOS
% ============================================================================
\section{Exercícios}

\begin{frame}{Exercícios - Parte 1}
\begin{block}{Questão 1: Crescimento Populacional}
Uma população segue o modelo logístico com $r = 0.6$ dia$^{-1}$ e $K = 5000$.
\begin{enumerate}
    \item Se $N(0) = 500$, calcule $N(10)$
    \item Em que tempo a população atinge 90\% de $K$?
    \item Qual a taxa máxima de crescimento ($dN/dt$)? Em que densidade?
\end{enumerate}
\end{block}

\begin{block}{Questão 2: Efeito Allee}
Considere o modelo com efeito Allee: $\frac{dN}{dt} = rN(1 - N/K)(N/A - 1)$, com $r = 0.5$, $K = 1000$, $A = 100$.
\begin{enumerate}
    \item Encontre os pontos de equilíbrio
    \item Analise a estabilidade de cada equilíbrio
    \item Simule em Python para $N(0) = 50$ e $N(0) = 150$
    \item Qual a implicação para conservação?
\end{enumerate}
\end{block}
\end{frame}

\begin{frame}{Exercícios - Parte 2}
\begin{block}{Questão 3: Modelo de Ricker}
Implemente o modelo de Ricker: $N_{t+1} = N_t e^{r(1-N_t/K)}$
\begin{enumerate}
    \item Simule para $r = 1.5, 2.0, 2.5, 3.0$ com $K = 100$
    \item Crie diagramas de bifurcação variando $r$ de 0 a 3
    \item Identifique regiões de estabilidade e caos
    \item Compare com o modelo de Beverton-Holt
\end{enumerate}
\end{block}

\begin{block}{Questão 4: Lotka-Volterra}
Para o sistema predador-presa com $\alpha = 1.0$, $\beta = 0.1$, $\delta = 0.75$, $\gamma = 1.5$:
\begin{enumerate}
    \item Calcule o ponto de equilíbrio não-trivial
    \item Calcule a matriz Jacobiana no equilíbrio
    \item Determine os autovalores. O equilíbrio é estável?
    \item Simule e plote séries temporais e retrato de fase
\end{enumerate}
\end{block}
\end{frame}

\begin{frame}{Exercícios - Parte 3}
\begin{block}{Questão 5: Modelo SIR}
Uma doença tem $\beta = 0.4$ dia$^{-1}$ e $\gamma = 0.1$ dia$^{-1}$ em uma população de $N = 10000$.
\begin{enumerate}
    \item Calcule $R_0$
    \item Se inicialmente $I(0) = 20$, simule 200 dias
    \item Qual o pico de infectados? Quando ocorre?
    \item Qual a taxa de ataque final (proporção que adoece)?
    \item Que fração deve ser vacinada para prevenir epidemia?
\end{enumerate}
\end{block}

\begin{alertblock}{Projeto: Competição e Coexistência}
Implemente o modelo de competição Lotka-Volterra:
\begin{itemize}
    \item Simule os 4 cenários (sp1 vence, sp2 vence, coexistência, bistabilidade)
    \item Crie retratos de fase com isoclinas
    \item Adicione estocasticidade demográfica
    \item Explore efeitos de flutuações ambientais
\end{itemize}
\end{alertblock}
\end{frame}

% ============================================================================
% SEÇÃO 7: REFERÊNCIAS
% ============================================================================
\section{Referências}

\begin{frame}{Referências Principais}
\begin{thebibliography}{99}
\bibitem{murray} Murray J.D. (2002). \textit{Mathematical Biology I: An Introduction}, 3rd ed. Springer.

\bibitem{edelstein} Edelstein-Keshet L. (2005). \textit{Mathematical Models in Biology}. SIAM.

\bibitem{otto} Otto S.P., Day T. (2007). \textit{A Biologist's Guide to Mathematical Modeling in Ecology and Evolution}. Princeton University Press.

\bibitem{keeling} Keeling M.J., Rohani P. (2008). \textit{Modeling Infectious Diseases in Humans and Animals}. Princeton University Press.

\bibitem{anderson} Anderson R.M., May R.M. (1991). \textit{Infectious Diseases of Humans: Dynamics and Control}. Oxford University Press.
\end{thebibliography}
\end{frame}

\begin{frame}{Leitura Complementar}
\begin{block}{Ecologia Teórica}
\begin{itemize}
    \item Case T.J. (2000). \textit{An Illustrated Guide to Theoretical Ecology}
    \item Kot M. (2001). \textit{Elements of Mathematical Ecology}
    \item Tilman D. (1982). \textit{Resource Competition and Community Structure}
\end{itemize}
\end{block}

\begin{block}{Epidemiologia}
\begin{itemize}
    \item Diekmann O., Heesterbeek H. (2000). \textit{Mathematical Epidemiology of Infectious Diseases}
    \item Brauer F., Castillo-Chavez C. (2012). \textit{Mathematical Models in Population Biology and Epidemiology}
\end{itemize}
\end{block}

\begin{block}{Recursos Computacionais}
\begin{itemize}
    \item \textbf{EpiModel:} \url{https://epimodel.org} (R package para epidemiologia)
    \item \textbf{PyDSTool:} \url{https://pydstool.github.io} (Python para sistemas dinâmicos)
    \item \textbf{COPASI:} \url{http://copasi.org} (simulação de modelos)
\end{itemize}
\end{block}
\end{frame}

% ============================================================================
% SLIDE FINAL
% ============================================================================
\begin{frame}
\begin{center}
{\Huge Obrigado!}

\vspace{1cm}

{\Large Capítulo 10: Sistemas Populacionais}

\vspace{1cm}

\textbf{Próximo Capítulo:}\\
Tópicos Avançados em Biologia de Sistemas

\vspace{1cm}

\textit{Dúvidas? Perguntas?}
\end{center}
\end{frame}

\end{document}
